\chapter{Introduction}
\label{chap:introduction}

% One sentence per line keeps diffs readable and makes a sentence easy to move.
Lorem ipsum dolor sit amet, consectetur adipiscing elit, sed do eiusmod tempor incididunt ut labore et dolore magna aliqua.
A citation looks like this~\cite{example_JournalArticle_2023}, and a citation with the author in the sentence reads \textcite{example_ConferencePaper_2024} instead.
Acronyms expand on first use, so an \ac{llm} prints in full here and as \ac{llm} from the second use on.
Cross-references go through cleveref, which supplies the word: \Cref{chap:background} is a chapter, \Cref{sec:method-protocol} a section, \Cref{tab:intro-example} a table.

\section{Motivation}
\label{sec:motivation}

Ut enim ad minim veniam, quis nostrud exercitation ullamco laboris nisi ut aliquip ex ea commodo consequat.
Duis aute irure dolor in reprehenderit in voluptate velit esse cillum dolore eu fugiat nulla pariatur.

\begin{margintable}
  \centering
  \caption[A table in the margin]{A margin table. It sits in the 41.4\,mm margin column, so keep it to three or four narrow columns.}
  \label{tab:intro-example}
  \small
  \begin{tabular}{@{}lrr@{}}
    \toprule
    Variant & F1 & Rec \\
    \midrule
    baseline & 0.501 & 0.206 \\
    proposed & 0.909 & 0.847 \\
    \bottomrule
  \end{tabular}
\end{margintable}

Excepteur sint occaecat cupidatat non proident, sunt in culpa qui officia deserunt mollit anim id est laborum.
Sed ut perspiciatis unde omnis iste natus error sit voluptatem accusantium doloremque laudantium.\sidenote{A sidenote lands in the margin beside the line that calls it, numbered like a footnote.}
Totam rem aperiam, eaque ipsa quae ab illo inventore veritatis et quasi architecto beatae vitae dicta sunt explicabo.

\section{Research Question}
\label{sec:research-question}

Nemo enim ipsam voluptatem quia voluptas sit aspernatur aut odit aut fugit, sed quia consequuntur magni dolores eos qui ratione voluptatem sequi nesciunt.

\begin{quote}
  How does one phrase a research question so that an experiment can answer it?
\end{quote}

Neque porro quisquam est, qui dolorem ipsum quia dolor sit amet, consectetur, adipisci velit.

\section{Contributions}
\label{sec:contributions}

\begin{enumerate}
  \item \textbf{The first contribution.} Sed quia non numquam eius modi tempora incidunt ut labore et dolore magnam aliquam quaerat voluptatem (\Cref{chap:methodology}).
  \item \textbf{The second contribution.} Ut enim ad minima veniam, quis nostrum exercitationem ullam corporis suscipit laboriosam (\Cref{chap:results}).
  \item \textbf{The third contribution.} Quis autem vel eum iure reprehenderit qui in ea voluptate velit esse quam nihil molestiae consequatur (\Cref{app:first}).
\end{enumerate}

\section{Thesis Outline}
\label{sec:thesis-outline}

\Cref{chap:background} reviews the literature, \Cref{chap:dataset} describes the corpus, \Cref{chap:methodology} defines the task and the evaluation protocol, \Cref{chap:results} reports the experiments, \Cref{chap:discussion} interprets them, and \Cref{chap:conclusion} closes with the findings and the open questions.
